
\clearpage
\VerdierPhysicalPageAnchor{0129}%
\begingroup
\small
\setlength{\emergencystretch}{1.5em}%
\setlength{\abovedisplayskip}{3pt plus 2pt minus 2pt}%
\setlength{\belowdisplayskip}{3pt plus 2pt minus 2pt}%

\VerdierSemanticAnchor{ch2-2-1-1-sm3-conditions}
\VerdierSemanticAnchor{ch2-2-1-1-sm3-i}
\noindent\textit{i)} 存在 $S$ 中源为 $Y$ 的态射 $s$，使 $sf=sg$。

\VerdierSemanticAnchor{ch2-2-1-1-sm3-ii}
\noindent\textit{ii)} 存在 $S$ 中靶为 $X$ 的态射 $t$，使 $ft=gt$。

\VerdierSemanticAnchor{ch2-2-1-1-saturated}
\noindent 一个乘法系称为\emph{饱和的}，如果它具有下列性质：

\VerdierSemanticAnchor{ch2-2-1-1-sm4}
\VerdierSemanticAnchor{ch2-2-1-1-saturated-sm4}
\smallskip
\noindent\textbf{SM4)} 如果 $f,g,$ 与 $h$ 是三个可复合态射，并且 $fg$ 与 $gh$
属于 $S$，则 $g$ 属于 $S$。

\VerdierSemanticAnchor{ch2-2-1-2}
\medskip
\noindent\textbf{2.1.2.}\quad
\VerdierSemanticAnchor{ch2-2-1-2-compatible-triangulation}
\VerdierSemanticAnchor{ch2-2-1-2-compatible}
设 $\mathcal D$ 为三角范畴。$\mathcal D$ 的乘法系 $S$ 称为
\emph{与三角结构相容}，如果它具有下列性质：

\VerdierSemanticAnchor{ch2-2-1-2-sm5}
\smallskip
\noindent\textbf{SM5)} 对于 $S$ 的每个元素 $s$ 与每个整数
$n\in\mathbb Z$，态射 $s[n]$ 属于 $S$。换言之，$S$
\emph{在平移自同构下稳定}。

\VerdierSemanticAnchor{ch2-2-1-2-sm6}
\smallskip
\noindent\textbf{SM6)} 对于任意一对特异三角：

\VerdierSemanticAnchor{ch2-2-1-2-sm6-triangle-pair}
\[
  (X,Y,Z,u,v,w), \qquad (X',Y',Z',u',v',w')
\]
\noindent 以及任意交换图：

\VerdierSemanticAnchor{ch2-2-1-2-sm6-square}
\begin{center}
  \begin{tikzcd}[column sep=4.2em,row sep=2.5em]
    X \arrow[r,"u"] \arrow[d,"s"'] & Y \arrow[d,"s'"] \\
    X' \arrow[r,"u'"'] & Y'
  \end{tikzcd}
  \quad ,
\end{center}
\VerdierSemanticAnchor{ch2-2-1-2-sm6-triangle-morphism}
\VerdierSemanticAnchor{ch2-2-1-2-sm6-completion}
\noindent 其中 $s$ 与 $s'$ 都属于 $S$，存在属于 $S$ 的态射
$s''\colon Z\to Z'$，使 $(s,s',s'')$ 是三角的态射。

\VerdierSemanticAnchor{ch2-2-1-3}
\medskip
\noindent\textbf{注 2.1.3.}\quad
\VerdierSemanticAnchor{ch2-2-1-3-item1}
\VerdierSemanticAnchor{ch2-2-1-3-1}
\textbf{1)} 设 $\mathcal D$ 为三角范畴，$S$ 为 $\mathcal D$ 中与三角结构
相容的乘法系，$\mathcal D^\circ$ 为 $\mathcal D$ 的反三角范畴。把 $S$ 看作
$\mathcal D^\circ$ 的态射集时，它是 $\mathcal D^\circ$ 的乘法系，并与
$\mathcal D^\circ$ 的三角结构相容 (1.1.7)。

\VerdierSemanticAnchor{ch2-2-1-3-item2}
\VerdierSemanticAnchor{ch2-2-1-3-2}
\smallskip
\noindent\textbf{2)} 设 $\mathcal D$ 为三角范畴，$S$ 为 $\mathcal D$ 中具有
性质 SM1)、SM5)、SM6) 的态射集。那么 $S$ 也具有性质 SM2)。
\VerdierSemanticAnchor{ch2-2-1-3-2-sm2}
事实上，只需证明任意图式：

\VerdierSemanticAnchor{ch2-2-1-3-item2-input-boundary}
\VerdierSemanticAnchor{ch2-2-1-3-2-input-boundary}
\begin{center}
  \begin{tikzcd}[column sep=4.0em,row sep=2.6em]
    & X' \arrow[d,"s\in S"] \\
    Z \arrow[r,"f"] & X
  \end{tikzcd}
\end{center}
\endgroup

